\documentclass[a4paper,11pt]{article}
\usepackage{exam-style-341}
\examinehead{341 农业综合知识三 · 程序设计}{模拟试卷（二）参考答案}
\begin{document}
\answercover{341 农业综合知识三 · 程序设计}{模拟试卷（二）参考答案}

\answerpart{一、单项选择题（每题 2 分，共 20 分）}
\begin{qgroup}
\q A。
\begin{analysis}
\code{a / b} 的两个操作数都是 \code{int}，按整数除法得 2；再与 \code{float} 型 \code{x} 相乘时，2 被转换为 \code{float}，结果为 6.0。若误按 \code{a * x / b} 计算才会得到 7.5。
\end{analysis}

\q A。
\begin{analysis}
\code{i++} 是先取 \code{i} 的原值再自增：表达式 \code{i++} 的值为 3，运算后 \code{i} 变为 4，故 \code{j = 3 + 2 = 5}，输出 \code{4 5}。
\end{analysis}

\q B。
\begin{analysis}
\code{do-while} 先执行循环体再判断条件。循环体依次使 \code{s} 增加 1、2、3、4，即 \code{s = 1 + 2 + 3 + 4 = 10}；此时 \code{i = 4}，条件 \code{i < 4} 不成立，循环结束。
\end{analysis}

\q A。
\begin{analysis}
循环取的是副对角线元素：\code{a[0][2] = 3}、\code{a[1][1] = 5}、\code{a[2][0] = 7}，累加得 15。
\end{analysis}

\q A。
\begin{analysis}
\code{p = s + 1} 指向字符 \code{'b'}，故 \code{*p} 输出 \code{b}；\code{p + 2} 指向 \code{s[3]}，即字符串 \code{"def"}，用 \code{\%s} 输出为 \code{def}。整个输出是 \code{bdef}。
\end{analysis}

\q B。
\begin{analysis}
按定义展开：\code{f(1) = 1}，\code{f(0) = 1}（满足 \code{n <= 1}），\code{f(2) = f(1) + f(0) = 2}，\code{f(3) = f(2) + f(1) = 3}，\code{f(4) = f(3) + f(2) = 5}，\code{f(5) = f(4) + f(3) = 8}。
\end{analysis}

\q A。
\begin{analysis}
\code{*p = *q} 把 \code{b} 的值赋给 \code{a}，\code{a = 20}；\code{q = p} 使 \code{q} 也指向 \code{a}；\code{*q = 30} 修改的是 \code{a}，故 \code{a = 30}。全程没有语句修改 \code{b}，\code{b} 仍为 20，输出 \code{30 20}。
\end{analysis}

\q C。
\begin{analysis}
\code{a.next = \&b} 使 \code{a} 的指针成员指向结构体变量 \code{b}，因此 \code{a.next->x} 等于 \code{b.x = 2}，而 \code{a.x = 1}，两者之和为 3。
\end{analysis}

\q A。
\begin{analysis}
\code{fopen} 打开失败时返回空指针 \code{NULL}，程序必须判断返回值，否则后续读写会导致严重错误，A 正确。文件一旦被 \code{fclose} 关闭，该文件指针就不再指向有效的文件流，继续用它读写属于非法操作，B 错误；\code{fgetc} 一次只读入一个字符，读字符串要用 \code{fgets} 或 \code{fscanf}，C 错误；以 \code{"w"} 方式打开已存在的文件会清空其内容，D 错误。
\end{analysis}

\q A。
\begin{analysis}
\code{atoi} 把数字字符串转换为 \code{int}；\code{atof} 转换为 \code{double}；\code{itoa} 不是标准 C 库函数；\code{sprintf} 是把格式化的数据写入字符串（方向相反）。
\end{analysis}
\end{qgroup}

\answerpart{二、填空题（每题 2 分，共 10 分）}
\begin{qgroup}
\q 1。
\begin{analysis}
条件运算符是右结合的，该表达式等价于 \code{x > 0 ? (x > 5 ? 1 : 2) : 3}。\code{x = 10}，\code{x > 0} 成立，再判断 \code{x > 5} 也成立，故取内层表达式的值 1。
\end{analysis}

\q \code{n != 0}。
\begin{analysis}
辗转相除法的结束条件是除数 \code{n} 变为 0，此时被除数 \code{m} 就是最大公约数，故循环条件为 \code{n != 0}。若写成 \code{r != 0}，第一次判断时 \code{r} 尚未赋值，属于使用未初始化变量。
\end{analysis}

\q \code{'\textbackslash 0'}；10。
\begin{analysis}
用字符串常量 \code{"abc"} 初始化字符数组时，会额外存放一个字符串结束标志 \code{'\textbackslash 0'}，即 \code{s[3] = '\textbackslash 0'}；而数组 \code{s} 的大小由声明决定，\code{sizeof(s) = 10}，未显式赋值的元素自动为 0。
\end{analysis}

\q 4；3。
\begin{analysis}
\code{a + 3} 指向 \code{a[3]}，所以 \code{*p = 4}；两个指向同一数组的指针相减，结果是它们之间相差的元素个数，即 \code{(a + 3) - a = 3}。
\end{analysis}

\q \code{fclose(fp)}。
\begin{analysis}
用 \code{"w"} 方式写入的数据先存放在缓冲区中，必须先调用 \code{fclose} 关闭文件，把缓冲区内容写入磁盘，然后才能用 \code{"r"} 方式重新打开并读到正确数据。
\end{analysis}
\end{qgroup}

\answerpart{三、程序阅读题（每题 5 分，共 10 分）}
\begin{qgroup}
\q 运行结果为：\code{10 13}
\begin{analysis}
\code{static int s = 0;} 是静态局部变量，只在程序开始执行时初始化一次，函数返回后它的值仍然保留。
\end{analysis}
\begin{analysis}
第一次调用 \code{f(4)}：\code{s = 0 + 4 = 4} $\rightarrow$ \code{s = 4 + 3 = 7} $\rightarrow$ \code{s = 7 + 2 = 9} $\rightarrow$ \code{s = 9 + 1 = 10}，最后调用 \code{f(0)} 直接返回 \code{s}，所以 \code{a = 10}。
\end{analysis}
\begin{analysis}
第二次调用 \code{f(2)} 时 \code{s} 仍为 10（static 变量未被重新初始化）：\code{s = 10 + 2 = 12} $\rightarrow$ \code{s = 12 + 1 = 13}，\code{f(0)} 返回 13，所以 \code{b = 13}。故输出 \code{10 13}。
\end{analysis}
\begin{analysis}
对比：若把 \code{s} 写成普通的自动变量（去掉 \code{static}），则 \code{f(4) = 10}、\code{f(2) = 3}，输出将变为 \code{10 3}。
\end{analysis}

\q 三处错误及其改正：
\begin{analysis}
第 8 行：循环条件 \code{i <= N} 会访问 \code{a[N]}（即 \code{a[5]}），造成数组下标越界，应改为 \code{i < N}。
\end{analysis}
\begin{analysis}
第 10 行：\code{avg = sum / N;} 中 \code{sum} 与 \code{N} 都是整型，进行的是整数除法，小数部分被截去，应改为 \code{avg = (float)sum / N;}（也可写成 \code{sum / 5.0} 或 \code{(float)sum / 5}）。
\end{analysis}
\begin{analysis}
第 11 行：\code{printf} 语句末尾缺少分号，应在 \code{)} 之后补上 \code{;}。
\end{analysis}
\begin{analysis}
改正后的完整程序：
\end{analysis}
\begin{lstlisting}[language={[custom]C}]
#include <stdio.h>
#define N 5
int main(void)
{
    int a[N] = {1, 2, 3, 4, 6};
    int i, sum = 0;
    float avg;
    for (i = 0; i < N; i++)
        sum += a[i];
    avg = (float)sum / N;
    printf("sum=%d, avg=%f\n", sum, avg);
    return 0;
}
\end{lstlisting}
\begin{analysis}
改正后数组元素之和为 $1+2+3+4+6=16$，平均值为 $16/5=3.2$，输出 \code{sum=16, avg=3.200000}。若保持原来的整数除法，\code{avg} 将得到 3.000000，与正确结果不符。
\end{analysis}
\end{qgroup}

\answerpart{四、程序设计题（每题 5 分，共 10 分）}
\begin{qgroup}
\q 参考实现：
\begin{lstlisting}[language={[custom]C}]
#include <stdio.h>
#include <string.h>
void reverse(char *s) {
    char *p = s, *q = s + strlen(s), t;
    if (p == q) return;            /* 空串，无需逆序 */
    q--;                           /* q 指向最后一个有效字符 */
    while (p < q) {
        t = *p; *p = *q; *q = t;    /* 交换首尾两个字符 */
        p++;
        q--;
    }
}
int main(void) {
    char str[100];
    printf("请输入一个字符串：");
    scanf("%s", str);
    reverse(str);
    printf("%s\n", str);
    return 0;
}
\end{lstlisting}
\begin{analysis}
评分要点：正确定义并使用两个指针（或两个下标）分别从首尾向中间移动（2 分）；正确交换 \code{*p} 与 \code{*q} 并保证循环在 \code{p < q} 时进行，不重复交换、不越界（2 分）；主函数输入输出与函数调用正确（1 分）。
\end{analysis}
\begin{analysis}
说明：\code{q = s + strlen(s) - 1} 与本题写法等价；若字符串为空串，\code{strlen(s)} 为 0，必须先判断再 \code{q--}，否则指针会指向数组之前，属于非法访问。另外 \code{scanf("\%s", str)} 遇到空格即结束，若要逆序整行（含空格）应改用 \code{fgets}。
\end{analysis}

\q 参考实现：
\begin{lstlisting}[language={[custom]C}]
#include <stdio.h>
#define N 5
struct student {
    int num;
    char name[20];
    float score[3];
    float avg;
};
int main(void) {
    struct student s[N];
    int i, j, maxi = 0;
    float sum;
    for (i = 0; i < N; i++) {              /* 读入并求平均分 */
        scanf("%d %s", &s[i].num, s[i].name);
        sum = 0.0;
        for (j = 0; j < 3; j++) {
            scanf("%f", &s[i].score[j]);
            sum += s[i].score[j];
        }
        s[i].avg = sum / 3;
    }
    for (i = 0; i < N; i++)                /* 输出每人平均分 */
        printf("%d %s %.2f\n", s[i].num, s[i].name, s[i].avg);
    for (i = 1; i < N; i++)                /* 找平均分最高者 */
        if (s[i].avg > s[maxi].avg) maxi = i;
    printf("平均分最高的学生：%s，平均分 %.2f\n", s[maxi].name, s[maxi].avg);
    return 0;
}
\end{lstlisting}
\begin{analysis}
评分要点：正确定义含学号、姓名、3 门成绩和平均分的结构体类型（1 分）；用结构体数组和嵌套循环完成输入与累加，平均分用 \code{sum / 3} 按实数计算（2 分）；用“打擂台”法求平均分最高者，注意 \code{maxi} 初始化为 0、比较从下标 1 开始（1 分）；正确输出姓名与平均分（1 分）。
\end{analysis}
\begin{analysis}
易错点：\code{scanf} 中数组名 \code{s[i].name} 本身就是地址，不能再加 \code{\&}；成绩是 \code{float} 型，输入必须用 \code{"\%f"}；比较平均分时不要写成字符串比较，应比较 \code{s[i].avg}。
\end{analysis}
\end{qgroup}

\end{document}
